% Local Validity and Contextual Holonomy -- v3 (full rewrite)
% Working draft -- pending adversarial review.
% Discipline: every claim is (a) a cited theorem of papers A-D, (b) a named machine
% verification in manuflog/contextuality-obstructions, or (c) explicitly labeled
% Conjecture/Open. No claim of readiness is made in this file.
\documentclass[11pt]{article}
\usepackage{amsmath,amssymb,amsthm,geometry,booktabs,hyperref,enumitem}
\geometry{margin=1.1in}
\newtheorem{theorem}{Theorem}
\newtheorem{proposition}{Proposition}
\newtheorem{corollary}{Corollary}
\newtheorem{conjecture}{Conjecture}
\newtheorem{openproblem}{Open Problem}
\theoremstyle{definition}
\newtheorem{principle}{Principle}
\theoremstyle{remark}
\newtheorem{remark}{Remark}
\newcommand{\U}{\mathrm{U}}
\newcommand{\PU}{\mathrm{PU}}
\newcommand{\CF}{\mathrm{CF}}
\newcommand{\tr}{\operatorname{tr}}
\title{Local Validity and Contextual Holonomy:\\ A Structural Formulation of the Copenhagen Principle\\[2pt] \large (v3 --- complete rewrite)}
\author{Manuel Flores Gordillo\\ \small Independent Researcher \quad ORCID 0009-0006-8246-0343\\ \small \texttt{manuel.flores@columbia.edu}}
\date{July 2026 --- draft v3, pending adversarial review}
\begin{document}
\maketitle

\begin{abstract}
We propose that the Copenhagen viewpoint, stripped of its rhetoric, is a set of
structural statements about one geometric object: the transformation groupoid
$\PU(n)\ltimes X$ of measurement contexts (maximal abelian subalgebras) with its
canonical $\U(1)$ $2$-cocycle. In this formulation: \emph{local validity} --- the exact
classicality of every single context --- is a theorem with an explicit witness;
\emph{complementarity} is the nonvanishing holonomy of the cocycle along closed context
loops, which detects exactly the all-versus-nothing sector of contextuality (an equivalence,
now proved), with the
achievable state-independent values classified (the value classification --- even/odd dichotomy and
fixed value $d/2$ --- is due to Abramsky, Cercelescu \& Constantin, FSCD 2024, and is recast here in
certificate language) and the state sector governed by a facet
geometry that is now completely enumerated in the certified $d=4$ example ($23{,}256$
facets in $61$ symmetry classes) and classified in closed form at $d=2$;
\emph{collapse} is pinned by statistics, repeatability, and covariance to a
one-parameter family per degenerate block whose unique efficient member is the
L\"uders rule --- with a classical twin theorem identifying the same axioms' classical
solution as Bayesian conditioning; \emph{observer-relativity of facts} and the \emph{movable cut}
are theorems of an enlarged observer-context category; and the Frauchiger--Renner
argument localizes to an ordinary state-dependent violation (contextual fraction exactly
$1/6$) carried by a single CHSH facet, with the loop holonomy of the observer switch
provably trivial. Two of these signatures have been measured on superconducting hardware.
Each tenet is either a theorem of the companion papers, a machine-verified certificate in
the public suite, or is explicitly labeled a conjecture. Version v1 of this note
contained four errors of substance; \S\ref{sec:changed} states them and their repairs.
\end{abstract}

\section{Introduction: the principle, restated}\label{sec:intro}
Textbook presentations of the Copenhagen interpretation bundle together several
commitments: that classical concepts apply exactly within a single experimental
arrangement; that arrangements exclude one another (complementarity); that measurement
updates the state by projection; that facts are relative to the experimental context,
including the observer; and that the boundary between quantum system and classical
apparatus --- the cut --- can be moved without empirical consequence. Critics have long
held that these are attitudes rather than assertions. This note argues the opposite:
each commitment is a precise structural statement about a single mathematical object,
and each is now either proved, machine-certified in a public verification suite, or
stated as an explicit conjecture with its current evidence.

The object is the transformation groupoid
\[
\mathcal G \;=\; \PU(n)\ltimes X,
\]
where $X$ is the space of maximal abelian $*$-subalgebras (MASAs, i.e.\ measurement
contexts) of $M_n(\mathbb C)$ and the projective unitary group acts by conjugation,
together with the canonical $\U(1)$ $2$-cocycle obtained from the central extension
$\U(1)\to\U(n)\to\PU(n)$. The formulation is then:

\begin{principle}[Copenhagen, structural form]\label{prin}
\emph{(i) Local validity.} Every single context is exactly classical: the restriction of
the canonical cocycle to any one context trivializes, with an explicit trivializing
cochain. \emph{(ii) Complementarity.} The cocycle does not trivialize globally; its
holonomy along closed context loops is a quantized obstruction, and this obstruction is
the all-versus-nothing sector of contextuality, with the state sector governed by an
enumerable facet geometry. \emph{(iii) Collapse.} Statistics, repeatability, and covariance under the measured
observable's commutant pin the
measurement update to a one-parameter family per degenerate block, with the L\"uders
rule its unique efficient member; classically the same axioms force Bayesian
conditioning up to the same per-cell parameter. \emph{(iv) Relative facts.} Value assignments
exist only per context; the obstruction persists under enlargement of the observer.
\emph{(v) The movable cut.} The placement of the system/apparatus boundary is a choice
of chart; section-fixed comparisons of differently-cut descriptions agree.
\end{principle}

Throughout, ``verified'' means: reproduced by the named script in the companion
repository \cite{repo}, with the expected one-line output pinned in
\texttt{verification/INDEX.md}. The supporting mathematics is developed in the companion
papers \cite{paperA,paperB,paperC,paperD}; this note is the synthesis; proofs
live in the companion papers and the verification supplement, with the source of every
numbered claim listed in Appendix~A.

\section{What changed since v1}\label{sec:changed}
Version v1 (February 2026) contained four errors of substance, found by adversarial
review and machine checking. We state them because the credibility of a structural
program rests on its error ledger.
\begin{enumerate}[leftmargin=1.4em]
\item \emph{Wrong groupoid.} v1 built the cocycle on $\U(n)\ltimes X$. The lifting class
there is vacuous (Theorem~1 of \cite{paperA}; \texttt{thmG\_general.py}); the canonical
class lives on $\PU(n)\ltimes X$, where it has order exactly $n$
\cite[Thm.~8]{paperA}.
\item \emph{Overclaim ``encodes''.} v1 asserted that holonomy \emph{encodes}
contextuality. False in both directions: at $d=3$ the class is nontrivial of order $3$
while the achievable state-independent all-versus-nothing value is $0$
(\texttt{d3\_gap\_certificate.py}), and KCBS contextuality is invisible to the class
(\texttt{kcbs\_converse.py}). The correct statement is \emph{detects the AvN sector}
(Theorem~\ref{thm:avn} below).
\item \emph{Level-unspecified Wigner--friend holonomy.} v1 attributed $-i$ to ``the''
holonomy of the friend loop. The projective (ray-level) holonomy is $+1$; $-i$ is the
determinant-\emph{section} holonomy, with section-independent square $-1$
(\texttt{wf\_loop\_holonomy.py}). Claims must pin the level.
\item \emph{Efficiency-restricted L\"uders uniqueness.} v1's uniqueness argument for the
update rule quantified only over efficient (single-Kraus) instruments. The
instrument-level theorem without that restriction is Theorem~\ref{thm:lueders} below.
\end{enumerate}

\paragraph{Corrections (2026-07-11, second external audit).} A follow-up adversarial pass found
further defects in the measurement-update layer, now repaired (the core obstruction-spectrum results
of \cite{paperA,paperB} are unaffected): (a) the collapse family's complete-positivity range was
misstated --- on a rank-$r$ block it is the affine interval $-1/(r^2-1)\le p\le 1$, not $[0,1]$
(Theorem~\ref{thm:lueders}); (b) the sharp-core counterexample $\mathrm{diag}(1,c)$ was invalid (its
second outcome is not repeatable) and is replaced by a qutrit example (Theorem~\ref{thm:sharp}); (c)
the classical twin's range is $-1/(m-1)\le t\le 1$ with a second deterministic member (the swap) at
$m=2$ (Theorem~\ref{thm:twin}); (d) the V37 evidence is a numerical stress test, not a
``machine-exact'' certificate; (e) the observer structure of \cite{paperD} is a \emph{category} with a
switch subgroupoid, not a groupoid; and (f) the hardware significance is statistical (shot-noise) under
the implemented witness model, conditional on compatibility/nondisturbance assumptions.

\section{Local validity is exact}\label{sec:local}
\begin{theorem}[Local classicality]\label{thm:local}
Let $C$ be any single Weyl context at local dimension $d$. The restriction of the
canonical cocycle to $C$ (the Weyl cocycle on the context's label group) is a
coboundary, with a trivializing cochain $\lambda$ valued in $\mu_d$ for \emph{even}
$d$ and in $\mu_{2d}$ for \emph{odd} $d$; the odd-$d$ bound is sharp (an explicit
$d=3$ group requires $\mu_{2d}$).
\end{theorem}
\begin{proof}[Proof sketch]
On a commuting label group $G$ the parity of the $\tau$-exponent $c(u,v)$ equals
$\langle u,v\rangle \bmod 2$, which vanishes for even $d$; hence $c=2c'$ with $c'$ a
symmetric $\mathbb Z_d$-valued $2$-cocycle. Symmetric classes form
$\mathrm{Ext}(G,\mathbb Z_d)$, additive over a cyclic decomposition of $G$ and
detected by the cyclic transgressions $\sum_{j<a} c'(jg,g)=\tfrac12 a(a-2)q(g)-dK$,
which are divisible by $a=\mathrm{ord}(g)$ in both parity cases; every cyclic class
therefore vanishes, and so does the total class. For odd $d$ the same transgression
argument at level $2d$ gives $\mu_{2d}$. Certified, with complete searches and the
$d=3$ sharpness witness, in \texttt{mu\_d\_valuation.py} (V46).
\end{proof}
An immediate corollary of the parity lemma: at even $d$ every commuting closed
context has \emph{even} $\tau$-level phase --- context products are always
$\omega$-powers, and integrality is automatic rather than a gauge choice. (An earlier
claim that half-steps $\tau\notin\mu_d$ are necessary at even $d$ traced to a label-layout
artifact in a verification harness and is retracted; see V44/V46 in the suite index.
Genuine $\tau$-steps are an odd-$d$ phenomenon. The \emph{state-sector}
$\tau$-necessity results of \cite{paperC} are a different statement and are
unaffected.) Explicit per-context trivializers: \texttt{local\_classicality.py} (V23),
for all rows and columns of the Peres--Mermin square at $d=2$, a $d=3$ sample family,
and the certified $d=4$ family; the same script exhibits the global contrast (no joint
$\lambda$ exists across contexts where the obstruction is present, cross-referencing
V10/V13/V33). This is tenet (i) of Principle~\ref{prin} in exact form: every chart of
the atlas is flat, with the flattening written down.

\section{The global twist: what holonomy detects}\label{sec:global}
\begin{theorem}[Canonical class; \cite{paperA}]\label{thm:class}
On $\U(n)\ltimes X$ the $\U(1)$ lifting class vanishes identically. After Morita
reduction, the canonical class on $\PU(n)\ltimes X$ has order exactly $n$.
\end{theorem}

\begin{theorem}[Detection equivalence]\label{thm:avn}
Let $F$ be a closed-triple Weyl context family at even $d$, with incidence matrix $A$
and context-phase vector $s\in\mathbb Z_d^{\mathrm{ctx}}$ (canonical: by the parity
lemma of Theorem~\ref{thm:local}, commuting context phases at even $d$ are
automatically $\omega$-powers). The following
are equivalent: (a) no noncontextual value assignment over $\mathbb Z_d$ exists; (b)
some left-kernel vector $u$ ($uA\equiv 0 \bmod d$) has $u\cdot s\equiv d/2 \pmod d$;
(c) some certificate carries odd commutator carry \cite[Thm.~4]{paperB}. The achievable
state-independent values are exactly $\{0,d/2\}$ for even $d$ and $\{0\}$ for odd $d$
\cite[Thm.~1]{paperB}.
\end{theorem}
\begin{proof}
(b)$\Rightarrow$(a): a solution $\gamma$ would give $u\cdot s\equiv (uA)\gamma\equiv
0$. (a)$\Rightarrow$(b): by Smith normal form over $\mathbb Z$, unsolvability of
$A\gamma\equiv s \pmod d$ produces, constructively, $u=\bigl(d/\gcd(g_i,d)\bigr)U_i$ with $uA\equiv 0$
and $u\cdot s\not\equiv 0 \pmod d$ (double annihilation over $\mathbb Z_d$); such a
$u$ is a certificate of the multiplicity covered by \cite[Rem.~3]{paperB}, so by the
Obstruction Spectrum Theorem \cite[Thm.~1]{paperB} its value satisfies $2(u\cdot
s)\equiv 0$, forcing $u\cdot s\equiv d/2$. (b)$\Leftrightarrow$(c) is
\cite[Thm.~4]{paperB}. Finally, no gauge hypothesis is needed: by the parity lemma of
Theorem~\ref{thm:local}, every commuting closed context at even $d$ has even
$\tau$-level phase, so $s=s^{(2d)}/2\in\mathbb Z_d^{\mathrm{ctx}}$ is canonical. (For
general closed-triple data --- commuting or not --- the $q$-parity identity
$s^{(2d)}\equiv A\bar q \pmod 2$, $\bar q_v=q(v)\bmod 2$, obtained by telescoping the
Weyl product law, supplies the constructive gauge; kernel pairings are
gauge-invariant.) Certified end-to-end, with random and pinned positives, in
\texttt{solvability\_equivalence.py} (V44, layout-corrected),
\texttt{tau\_gauge\_general.py} (V45, including the $d=6$ extension), and
\texttt{mu\_d\_valuation.py} (V46).
\end{proof}
What was, in v2, a machine-verified converse ($2{,}100$ families, zero mismatches) is
thus a corollary of \cite{paperB}'s proven theorems plus five lines of ring theory ---
the gap was assembly, not mathematics.
Both containments of the naive slogan fail, and strictly
(\texttt{d3\_gap\_certificate.py}, \texttt{kcbs\_converse.py}): the class does not imply
operational contextuality, and state-dependent contextuality can be class-invisible.
Complementarity, tenet (ii), is therefore \emph{located}: it is the AvN obstruction, no
more and no less at the level of the canonical class.

\subsection{The state sector: a facet geometry, now enumerated}\label{sec:state}
State-dependent contextuality is quantified by the contextual fraction $\CF$. The
following results replace v1's vague appeal to ``holonomy of the state'' with exact
geometry.

\begin{theorem}[Codeword-polytope classification, $d=2$; V31/V32]\label{thm:d2}
For closed-triple Weyl families at $d=2$, $\CF(F,\rho)=0$ if and only if the correlation
vector $c(\rho)$ lies in the convex hull of the codeword set $L(F)$; quantitatively
$\CF=\max\bigl(0,(S^*-1)/2\bigr)$, where $S^*$ is the optimal facet value, proved
end-to-end via a completion lemma and Abramsky--Barbosa--Mansfield duality.
\end{theorem}

\begin{proposition}[On-hull]\label{lem:onhull}
Every quantum behavior of the certified five-context $d=4$ family of
Theorem~\ref{thm:census} satisfies three families of exact linear
identities: per-context normalization, shared-observable marginal consistency, and
vanishing on spectrally forbidden joint outcomes ($24$ constraints). The affine space
they cut out has dimension exactly $33$, equal to the affine span of the $64$
deterministic behaviors, which satisfy the same identities. Hence quantum moment
vectors lie in the deterministic affine hull.
\end{proposition}

\begin{theorem}[Complete facet census, certified $d=4$ family; V43]\label{thm:census}
For the certified five-context $d=4$ family with one context deleted, the
noncontextual polytope has exactly $23{,}256$ facets in $61$ orbit classes under the
derived vertex-symmetry group of order $768$. By Proposition~\ref{lem:onhull}, $\CF(\rho)>0$ if and only if $\mu(\rho)$
violates one of the $23{,}256$ inequalities.
\end{theorem}
The census was derived twice by independent methods --- an adjacency decomposition on
the facet--ridge graph with every pivot certified in exact integer arithmetic, and a
direct exact enumeration by \textsc{Normaliz} --- which agree class-for-class
(\texttt{d4\_facet\_census.py}). Structural laws of this geometry are proved in
\cite{paperC}: the ghost--tower law $\gamma-s^*\equiv d/2$ for context deletions, and
the necessity of the $\tau$-level (half-step) data for optimal separation
(\texttt{d4\_moment\_facets.py}).

\begin{conjecture}[Scenario-paired classification]\label{conj:pair}
For every closed-triple Weyl family at even $d$, $\CF(F,\rho)=0$ iff $c(\rho)\in
\operatorname{conv}L(F)$ in the scenario-paired formulation of \cite{paperC}.
\end{conjecture}
The global-frame formulations of this statement are refuted (V28--V30); the pairing is
forced, not optional. The status is now sharply split by the choice of paired
coordinates (V48/V52). In the \emph{full-character} reading ($c(\rho)$ = all
per-context characters), the classification is a \emph{theorem}: the support lemma
confines outputs to the spectrally allowed set, the paired characters determine each
context's distribution by Fourier completeness, and an LP mass identity converts hull
membership into an exact mixture --- in this reading the statement is
near-definitional, which is why the $5{,}556$-record harness (V48; $d=4,8$,
census-arbitrated, zero mismatches) could not fail: it certifies implementations. The
substantive open form is the \emph{compressed} reading (the correlator-type
coordinates of \cite{paperC}, whose $d=2$ case is Theorem~\ref{thm:d2}'s actual
content). Its residual is \emph{not} an on-hull rank identity: the natural
generalization of Prop.~\ref{lem:onhull} to arbitrary families is \emph{false} (V52:
pinned families with affine-dimension gaps up to $74$ at $d=8$; the census family's
rank equality is special). The correct reduction (proven, V52): quantum behaviors lie
in the deterministic affine hull iff a finite per-family \emph{phase-coherence law}
holds --- characters with equal restriction to the section lattice carry a single
Weyl label and phase; the label half is proven (Smith-form annihilator plus closure),
and the phase half --- a cross-context $\tau$-holonomy congruence of the same species
as the $q$-parity identity --- is the named open lemma, exactly verified on all $36$
tested families at $d=4,6,8$.

\begin{theorem}[Tower holonomy, resolved by cases; V47]\label{conj:tower}
(i) \emph{Operator level:} every closed Weyl word has holonomy in $\mu_{2d}$, and
every value of $\mu_{2d}$ is attained (proved for even $d$; verified instances at odd
$d$). (ii) \emph{Determinant-section level} (canonical $\tau$-grid section): at even
$d$ all loop holonomies lie in $\mu_{2d}$; the attained group equals $\mu_{2d}$
exactly when $d\equiv 2 \pmod 4$, and is \emph{constant} $\mu_4$ along the $2$-adic
tower $d=2,4,8,16$.
\end{theorem}
The original conjecture's exhaustiveness clause is thus \emph{refuted}: the
coincidence $\mu_4=\mu_{2d}$ at the base $d=2$ does not propagate up the tower. At odd
$d$ the $\mu_{2d}$ bound is section-dependent, and both directions are now
\emph{proven} (V51): the DFT section violates it exactly when $d\equiv 3\pmod 4$
(Schur's determinant $\det F = i^{d(d-1)/2}$; pinned instance: the $d=3$ observer loop
has holonomy $e^{i\pi/6}$), while in the Weil-normalized section every loop holonomy
lies in $\mu_{2d}$ for all odd $d$ --- the Gauss phase cancels the $\mu_4$ part of
Schur's phase exactly ($\det F_W=(-1)^{\lfloor d/4\rfloor}$), $\det M_a$ is the Jacobi
symbol $(a|d)$ by Zolotarev, and the Weil-attained group is $\mu_6$ if $3\mid d$ and
$\mu_2$ otherwise, exhaustive only at $d=3$. Claims must pin the section --- the same
lesson as v1's error (3).

\section{Collapse: the admissible update rule, and its classical twin}\label{sec:collapse}
\begin{theorem}[Instrument-level L\"uders; V37]\label{thm:lueders}
Let $\{\mathcal I_j\}$ be an instrument for a projective measurement with degenerate
blocks, satisfying (i) statistics, (ii) repeatability, and (iii) covariance under the
commutant of the measured observable. Then on each degenerate block of rank $r_j$,
\[
\mathcal I_j(\rho)=p_j\,P_j\rho P_j+(1-p_j)\,\operatorname{tr}(P_j\rho)\,\frac{P_j}{r_j},
\qquad -\frac{1}{r_j^{2}-1}\le p_j\le 1,
\]
an \emph{affine} depolarizing family (the lower endpoint is forced by complete positivity; this is
\emph{not} only the convex segment $p_j\in[0,1]$). The linear solution space has exactly one parameter
per degenerate block (nullity $2$ at $(4,[2,2])$, verified also at $(6,[3,2,1])$), and the L\"uders
instrument $p_j=1$ is its unique efficient (single-Kraus) member. Covariance under the MASA alone is
strictly insufficient. Proof via Schur's lemma on block-covariant channels together with the Choi
bound; the V37 script is a numerical affine-dimension stress test, and the CP interval is certified
exactly in \texttt{lueders\_cp\_interval.py}.
\end{theorem}

\begin{theorem}[Sharp core; V40]\label{thm:sharp}
For a POVM instrument, statistics plus repeatability force every output of
$\mathcal I_j$ into the eigenvalue-$1$ subspace of its effect $E_j$, and these
subspaces are mutually orthogonal. A strictly unsharp effect (no $1$-eigenvector)
admits no repeatable instrument at all. The folk claim ``repeatable $\Rightarrow$
projective'' is nevertheless false as stated, but a valid witness must be at least
three-dimensional: the qutrit two-effect POVM $E_1=\mathrm{diag}(1,0,c)$,
$E_2=\mathrm{diag}(0,1,1-c)$, $0<c<1$, admits a repeatable covariant instrument on \emph{every}
outcome while $E_1$ is non-projective. (An earlier qubit example $\mathrm{diag}(1,c)$ was
withdrawn: its second effect $\mathrm{diag}(0,1-c)$ has no eigenvalue $1$, so that outcome is not
repeatable.)
\end{theorem}

\begin{theorem}[Classical twin; V40]\label{thm:twin}
The same axioms, interpreted for classical stochastic instruments on a finite outcome
partition, force the affine family $T_t=t\,I+(1-t)\,J/m$ per cell of size $m$, with
$-\tfrac{1}{m-1}\le t\le 1$ (nonnegativity), the nullities matching the quantum case
block-for-block. Its deterministic covariant members are the identity ($t=1$, Bayesian
conditioning) for every $m$ and, in addition, the swap ($t=-1$) when $m=2$; so Bayes is the unique
deterministic member only for $m\ge 3$.
\end{theorem}

\begin{remark}[Position]
Theorems~\ref{thm:lueders}--\ref{thm:twin} give the Leifer--Spekkens thesis a sharpened
form: \emph{conditional on} statistics, repeatability and full commutant covariance, the
admissible update rule is restricted to an affine depolarizing family whose unique single-Kraus
member is L\"uders, with Bayes' rule as the classical shadow of the same axioms. This is a
characterization of admissible instruments --- not a derivation that collapse occurs, nor a
resolution of the definite-outcome problem; rival interpretations can use the same formal result.
What is quantum is not the updating; it is the geometry (\S\ref{sec:global}) on which the updating
lives.
\end{remark}

\section{Observers, relative facts, and the cut}\label{sec:observer}
The observer enters through the observer-context category of \cite{paperD}, whose
objects are pairs (context, observer chain) and whose morphisms are switches within a fixed
enlargement (isomorphisms, forming a subgroupoid) together with refinement isometries into larger
Hilbert spaces (non-invertible); the presence of non-invertible refinements makes this a
\emph{category}, not a groupoid, and holonomy invariants are taken on the switch subgroupoid.

\begin{theorem}[Relative facts; \cite{paperD}]\label{thm:relative}
Value assignments exist per object and pull back monotonically along fibered morphisms;
the global obstruction persists under Wigner enlargement --- for the Peres--Mermin
square tensored with a friend ($d=8$), the context signs remain $(+1)^5(-1)^1$
(\texttt{observer\_groupoid.py}, certificate A).
\end{theorem}

\begin{theorem}[$\mathbb Z_2$ switching invariant; \cite{paperD}]\label{thm:z2}
The square of the switch-loop holonomy is a gauge-independent $\mathbb Z_2$ invariant
(proved via Weyl determinants $\pm1$): $-1$ on the mutually-unbiased-bases triangle,
$+1$ on within-context Weyl loops. The converse fails: a complementary $2$-cycle with
value $+1$ is pinned.
\end{theorem}

\begin{theorem}[Movable cut; \cite{paperD}]\label{thm:cut}
Section-fixed comparisons of descriptions differing by cut placement agree, with the
update rule of Theorem~\ref{thm:lueders} forced on both sides.
\end{theorem}

\begin{theorem}[Frauchiger--Renner localized; \cite{paperD}]\label{thm:fr}
In the observer-context category, the FR cycle carries trivial class: all switch-loop
phases are $+1$, fully gauge-invariantly. The paradox localizes to an ordinary
state-dependent violation: $\CF=1/6$ exactly, attained on the single CHSH facet
$\langle \bar Z Z\rangle-\langle \bar Z X\rangle+\langle \bar X X\rangle+\langle \bar X
Z\rangle\le 2$ at quantum value $7/3$ (\texttt{fr\_cycle.py}).
\end{theorem}
Theorem~\ref{thm:fr} \emph{refutes} v1's implicit suggestion (and \cite{paperD}'s own
initial conjecture) that the FR argument is a holonomy phenomenon: the switch holonomy
is trivial; what FR exercises is the state sector of \S\ref{sec:state}, with definite
coordinates and a definite price.

\paragraph{The Wigner--friend loop, level-pinned.}
For the loop $T\to T_x\to T_y\to T$ generated by $\mathbb 1\otimes H$, $\mathbb
1\otimes S$: the ray-level (Pancharatnam) holonomy is $+1$; the determinant-section
holonomy is exactly $-i$, a $\tau$-level invariant with section-independent square
$-1$ (\texttt{wf\_loop\_holonomy.py}).

\paragraph{Hardware.}
On \texttt{ibm\_fez} (job \texttt{d986q62f47jc73a7hm2g}): the interferometric loop
phase is $-2.93^\circ\pm1.55^\circ$, consistent with ray-level $+1$ and excluding $-i$
at ${\sim}56\sigma$; on the same device the Peres--Mermin witness attains
$S=4.6125\pm0.0173$, i.e.\ $35.4\sigma$ above the tight deterministic bound $4$
(\texttt{hardware/results/ibm\_fez\_holonomy\_20260709.json}). One machine, both
signatures, both sharp.

\section{The principle, assembled}\label{sec:assembled}
Principle~\ref{prin} now reads as a theorem schedule. (i) is
Theorem~\ref{thm:local}. (ii) is Theorems~\ref{thm:class}--\ref{thm:avn} for the
state-independent sector and Theorems~\ref{thm:d2}--\ref{thm:census} for the state
sector, with Conjecture~\ref{conj:pair} (now with $d=4,8$ evidence, V48) and the
case-resolved tower Theorem~\ref{conj:tower} marking the frontier. (iii)
is Theorems~\ref{thm:lueders}--\ref{thm:twin}. (iv) is
Theorems~\ref{thm:relative}--\ref{thm:z2} and the localization
Theorem~\ref{thm:fr}. (v) is Theorem~\ref{thm:cut}. Nothing here modifies quantum
predictions; the claim is that the Copenhagen commitments, correctly stated, are
structure rather than stance --- local flatness, one quantized global twist, forced
bookkeeping, chart-relative facts, and chart-choice for the cut.

\begin{openproblem}
(a) The cohomological home of class-invisible (KCBS-type) contextuality. Progress
(V49/V50): support-presheaf and scenario-complex cohomologies provably cannot capture
it (pinned no-gos); a $\U(1)$-holonomy-twisted spectral pairing on the
non-orthogonality graph characterizes the state sector exactly on odd cycles
($C_5,C_7,C_9$: $\CF=2\max(0,S-\alpha)$, all holonomy coordinates active, trivial
class Perron-maximal) and on the $\alpha=2$ antiholes; the once-mysterious
``$(-1)$-twist $=$ classical bound'' is the reflection identity $G\mapsto 2I-G$ and
holds exactly when $\alpha=2$. The Peres--Mermin milestone is achieved (V54): on the
$24$-ray configuration the same local system is Kochen--Specker--rigid with pure-torsion
holonomy, and the AvN $-1$ is its \emph{forced} holonomy invariant (transversal parity
theorem; $h(z_{\mathrm{obs}})=-1$), with flex dimension --- $0$ for PM, $2$ for KCBS
--- as the switch between the two layers. The general dictionary (flexes $=$ state
moduli; rigidity $+$ forced torsion $\Leftrightarrow$ AvN) remains conjectural. (b)
The general-$d$ facet census. The $d=4$ arithmetic profile is complete (V53): $18$
species, all facet functionals rational (the tier-2 ``$4\sqrt2$'' unit bounds were a
normalization artifact --- exact $\{0,\pm1\}$ lifts with integer bounds $\{6,7\}$
exist), $\tau$-level data necessary for all $61$ classes; the general-$d$ census
itself needs Normaliz-class compute (protocol pinned in V53). (c)
Conjecture~\ref{conj:pair} in the compressed reading: the $\Theta$-coherence
congruence of V52 (its label half is proven; the odd-$d$ Weil-section case of
Theorem~\ref{conj:tower} is now closed, V51).
\end{openproblem}

\appendix
\section{Claims-to-verification table}
\begin{center}\small
\begin{tabular}{lll}
\toprule
Claim & Source & Verification\\
\midrule
Thm.~\ref{thm:local} & this note & \texttt{local\_classicality.py} (V23); \texttt{mu\_d\_valuation.py} (V46)\\
Thm.~\ref{thm:class} & \cite{paperA} Thm.~1, 8 & \texttt{thmG\_general.py}\\
Thm.~\ref{thm:avn} & this note; \cite{paperB} Thm.~1, 4 & \texttt{solvability\_equivalence.py} (V44); \texttt{tau\_gauge\_general.py} (V45)\\
Thm.~\ref{conj:tower} & this note & \texttt{tower\_holonomy.py} (V47); \texttt{weil\_section\_theorem.py} (V51)\\
Conj.~\ref{conj:pair} (split status) & this note; \cite{paperC} & \texttt{scenario\_paired\_evidence.py} (V48); \texttt{paired\_classification\_theorem.py} (V52)\\
gaps (both strict) & \cite{paperC} & \texttt{d3\_gap\_certificate.py}, \texttt{kcbs\_converse.py}\\
Thm.~\ref{thm:d2} & \cite{paperC} Thm.~2 & \texttt{paired\_frame\_construction.py}, \texttt{ghost\_facet\_theorem.py}\\
Thm.~\ref{thm:census} & this note; \cite{paperC} & \texttt{d4\_facet\_census.py} (V43)\\
tower laws & \cite{paperC} & \texttt{d4\_ghost\_tower.py}, \texttt{d4\_moment\_facets.py}\\
Thm.~\ref{thm:lueders} & this note & \texttt{lueders\_instrument.py} (V37)\\
Thm.~\ref{thm:sharp}, \ref{thm:twin} & this note & \texttt{s4\_povm\_bayes.py} (V40)\\
Thm.~\ref{thm:relative}, \ref{thm:z2} & \cite{paperD} & \texttt{observer\_groupoid.py} (V38)\\
Thm.~\ref{thm:cut} & \cite{paperD} & V37 + \cite{paperD} \S3\\
Thm.~\ref{thm:fr} & \cite{paperD} & \texttt{fr\_cycle.py} (V39)\\
WF level split & \cite{paperC} & \texttt{wf\_loop\_holonomy.py}\\
Prop.~\ref{lem:onhull} & this note & \texttt{d4\_facet\_census.py} (exact-rank lemma)\\
S4/S5 proofs & v2 supplement & \texttt{note\_v2\_patch.md}\\
hardware & --- & \texttt{ibm\_fez} job \texttt{d986q62f47jc73a7hm2g}\\
\bottomrule
\end{tabular}
\end{center}

\begin{thebibliography}{9}
\bibitem{paperA} M.~Flores Gordillo, \emph{The Canonical Class of the MASA Context
Stack and the Peres--Mermin Obstruction}, working draft, 2026.
\bibitem{paperB} M.~Flores Gordillo, \emph{The Obstruction Spectrum of Weyl Context
Families and the 2-adic Tower}, working draft, 2026.
\bibitem{paperC} M.~Flores Gordillo, \emph{What Contextual Holonomy Detects}, working
draft, 2026.
\bibitem{paperD} M.~Flores Gordillo, \emph{The Observer-Context Category: Relative
Facts, a $\mathbb Z_2$ Switching Invariant, and the Movable Cut}, working draft v0.3,
2026.
\bibitem{repo} \texttt{github.com/manuflog/contextuality-obstructions}, verification
suite, 2026.
\end{thebibliography}
\end{document}
